The benchmark threshold control shows that operating-point effects extend beyond the Gaussian experiment, although their direction varies by classifier and condition. Under supplied labels, the complete framework's mean G-mean advantage over no intervention decreases from approximately $0.019$--$0.020$ at default rules to slightly negative differences after threshold selection. Under corrupted fitting labels, its advantages over no intervention decrease but remain positive on average.

Relative to SMOTE, the corrupted-label LR mean difference changes from $+0.0070$ to $-0.0029$, whereas the RF difference increases from $+0.0019$ to $+0.0131$. Threshold selection therefore does not uniformly eliminate or improve the complete framework's relative performance. None of the 12 selected-threshold comparisons is significant after Holm adjustment. These findings support classifier- and condition-dependent operating-point effects, without establishing equivalence or identifying augmentation-location selection as the sole source of any remaining benefit.

The control retains classifier parameters selected using default-rule validation. It assesses an additional threshold-selection step rather than joint classifier--threshold optimization. Inner-model score scales may transfer imperfectly to the final refitted model, and supplied validation labels are not verified ground truth. These limitations should qualify interpretation of the selected-threshold results.
